% !PRINTED 434
\weberpage{434}
%<<CONT
folglich ist das zweite Glied von (20) durch $p^2$ teilbar, und wenn
man $\gamma p^2$ dafür einsetzt, so folgt
\[
\alpha\delta - \beta\gamma = 1.
\]
Wir kommen also zu dem Resultat:

IV. \gesperrt{Die Invariantengleichung $F_p(v,u)$ hat immer dann
und nur dann mehrfache Wurzeln, wenn $u=j(\omega)$ ist, worin $\omega$
die Wurzel einer quadratischen Gleichung ist, deren Diskriminante
nicht durch $p$ teilbar ist, aber eine Lösung der Gleichung (18)
gestattet.}

Es gibt für ein gegebenes $p$ nur eine endliche Anzahl von
Werten $m$, die dieser Forderung entsprechen, da $m < 4p^2$ sein
muß. Die Werte $m=3$, $m=4$ sind darunter als spezielle
Fälle enthalten.

Da nach (18) $-m$ quadratischer Rest von $p$ ist, so ist $p$
durch eine Form der Diskriminante $-m$ darstellbar; $p^2$ ist aber
durch die Hauptform einer Diskriminante darstellbar. Die
Formenklasse, durch die $p$ darstellbar ist, ist also selbst entweder
die Hauptklasse oder sie gibt, mit sich selbst komponiert, die
Hauptklasse, ist also zweiseitig. Wir können das Theorem IV. also
auch folgendermaßen aussprechen:

V. \gesperrt{Die Wurzeln der Diskriminante der Invariantengleichung
$F_p(v,u)$ sind die singulären Invarianten $j(\omega)$, worin $\omega$
die Wurzel einer quadratischen Gleichung von negativer Diskriminante
$D$ ist, für die die Primzahl $p$ den Index 1 oder 2 hat.}

\begin{center}\rule{0.18\linewidth}{0.4pt}\end{center}
